\documentclass[11pt,a4paper]{article}
\usepackage[T1]{fontenc}
\usepackage[utf8]{inputenc}
\usepackage{lmodern}
\usepackage[a4paper,margin=27mm,headheight=15pt]{geometry}
\usepackage{amsmath,amssymb,amsthm,mathtools,bm,mathrsfs}
\usepackage{booktabs,array,longtable}
\usepackage{microtype}
\usepackage{xcolor}
\usepackage{xurl}
\usepackage{enumitem}
\usepackage{fancyhdr}
\usepackage{hyperref}
\hypersetup{colorlinks=true,linkcolor=blue!45!black,citecolor=blue!45!black,urlcolor=blue!45!black,
 pdftitle={Finite Resonance Certificates and Sharp Analytic Normalization for Hahn--Fuchsian Systems},
 pdfauthor={Research article prepared for Vladimir Reshetnikov},
 pdfsubject={Hahn series, Fuchsian systems, accumulation of exponents, normal forms and logarithmic obstructions}}
\urlstyle{same}
\numberwithin{equation}{section}
\newtheorem{theorem}{Theorem}[section]
\newtheorem{proposition}[theorem]{Proposition}
\newtheorem{lemma}[theorem]{Lemma}
\newtheorem{corollary}[theorem]{Corollary}
\theoremstyle{definition}
\newtheorem{definition}[theorem]{Definition}
\newtheorem{example}[theorem]{Example}
\theoremstyle{remark}
\newtheorem{remark}[theorem]{Remark}
% ed. (2026-09-29): unnumbered environment for ProveIt editorial notes; it
% leaves every numbered statement of the delivered article unchanged.
\newtheorem*{ednote}{Editorial note (ProveIt, 2026-09-29)}
\newcommand{\C}{\mathbb C}
\newcommand{\R}{\mathbb R}
\newcommand{\N}{\mathbb N}
\newcommand{\K}{\mathcal K}
\newcommand{\G}{\Gamma}
\newcommand{\Res}{\mathcal R}
\newcommand{\ad}{\operatorname{ad}}
\newcommand{\supp}{\operatorname{supp}}
\newcommand{\Mat}{\operatorname{Mat}}
\newcommand{\rank}{\operatorname{rank}}
\newcommand{\diag}{\operatorname{diag}}
\newcommand{\val}{\operatorname{v}}
\newcommand{\norm}[1]{\left\lVert #1\right\rVert}
\newcommand{\abs}[1]{\left\lvert #1\right\rvert}
\newcommand{\Aplus}{A_{+}}
\newcommand{\Hhat}{\widehat H}
\newcommand{\Ph}{\mathcal P}
\newcommand{\Th}{\mathcal T}
\newcommand{\Dg}{\mathcal D}
\newcommand{\Id}{\mathrm{id}}
\newcommand{\snap}{\texttt{04473354a0f3edff2d6c365caf26170ca6b88735}}
\DeclareMathOperator{\Sol}{Sol}
\setlist{itemsep=3pt,topsep=5pt,leftmargin=1.7em}
\setlength{\parskip}{3pt}
\setlength{\parindent}{1.2em}
\linespread{1.025}
\pagestyle{fancy}
\fancyhf{}
\fancyhead[L]{\small Hahn--Fuchsian systems}
\fancyhead[R]{\small Resonance and analytic normalization}
\fancyfoot[C]{\thepage}
\renewcommand{\headrulewidth}{0.3pt}
\setcounter{tocdepth}{1}
\begin{document}
% ed. (2026-09-29): no hyperref page anchor on the title page, whose page
% number 1 recurs after it (removes a duplicate-destination warning).
\hypersetup{pageanchor=false}
\begin{titlepage}
\thispagestyle{empty}
\vspace*{10mm}
\noindent{\large\scshape Transseries research}\par
\vspace{9mm}
\noindent{\LARGE\bfseries Finite Resonance Certificates\par}
\vspace{3mm}
\noindent{\LARGE\bfseries and Sharp Analytic Normalization\par}
\vspace{3mm}
\noindent{\LARGE\bfseries for Hahn--Fuchsian Systems\par}
\vspace{7mm}
\noindent{\large Accumulating exponents, unavoidable logarithms,\par
and a summability-threshold crossover}
\vspace{10mm}
\hrule
\vspace{6mm}
\noindent Research article prepared for Vladimir Reshetnikov\par
\noindent 29 September 2026
\vspace{8mm}
\begin{abstract}
We study linear systems $xY'=(A_0+\Aplus(x))Y$ whose positive real exponents
form a well-ordered set, without requiring local finiteness. For real-spectrum
$A_0$, we construct a normalized Hahn-series gauge and a finite resonant
normal form. A single nilpotent matrix determines the complete filtration by
logarithmic degree. Its entries depend on only finitely many original
coefficients, even when infinitely many exponents lie below the largest
resonance.

The central analytic result is a sharp universal criterion. Every absolutely
convergent, semigroup-supported coefficient series has an absolutely
convergent normalized gauge if and only if the nonresonant semigroup
exponents are uniformly separated from the eigenvalue differences. For a
well-ordered real semigroup, failure means accumulation from below at a
positive eigenvalue difference. A rank-one matrix perturbation gives
arbitrarily small counterexamples whenever this condition fails.

An explicit two-dimensional family has coefficient weights $n^{-p}$ and
exponents $2-1/n$. Its formal normalizer converges exactly when $p>2$,
although a renormalized analytic solution exists for every $p>1$ and realizes
all finite asymptotic prefixes. We derive its truncation crossover on the
scale $N\asymp\log(1/x)$, including a uniform Euler--Maclaurin error bound.
All general claims are proved conventionally. Reproducible finite symbolic
checks and numerical illustrations accompany the article. Classical
Fuchsian normal forms and the repository's earlier finite-obstruction work
are explicitly acknowledged; an exhaustive claim of historical priority is
not made.
\end{abstract}
\vfill
\noindent\small
\textbf{Scope.} Positive real Hahn exponents; a constant matrix with real
spectrum; one logarithmic indeterminate. This is not a solution of arbitrary
transseries differential equations and is not a Lean formalization.\par
\vspace{2mm}
\noindent\small
\textbf{Repository snapshot.}\quad\snap
\end{titlepage}
% ed. (2026-09-29): page anchors resume after the title page.
\hypersetup{pageanchor=true}
\tableofcontents
\newpage

\section{The research question and its place in ProveIt}
\label{sec:intro}

A formal transseries calculation and an analytic solution are different
objects. A well-ordered support makes coefficientwise algebra meaningful;
it does not bound the coefficients produced by division. This distinction
becomes decisive when exponents accumulate at a finite value. The familiar
Frobenius recurrence then meets a phenomenon absent from ordinary Taylor
series: infinitely many different exponents can approach one fixed
resonance.

The relevant ProveIt materials already develop asymptotic scales,
flatness, polynomial--logarithmic block calculus, Hahn-support lemmas, and
inversion methods. Their status documentation distinguishes actual Lean
results from human proofs and unformalized extensions
\cite{repo-index,repo-wellbased,repo-block}. A recently filed companion,
\emph{Finite Residue Obstructions and Logarithmic-Depth Promotion in Hahn
Transseries}, treats scalar differential operators by a finite obstruction
matrix under a uniform outer-exponent decrease \cite{repo-residue}.
Consequently, the mere assertion that resonances can have finite
obstruction data would not be a substantive new contribution to this
repository.

Here the principal question is instead the following:
\begin{quote}
When does a support-controlled formal Fuchsian normalization automatically
preserve absolute convergence, and what exactly goes wrong when it does
not?
\end{quote}
Two additional questions make the answer useful. Can all unavoidable
logarithms be certified from finite input data even for non-locally-finite
supports? Can a divergent normalizer nevertheless describe the finite
asymptotic prefixes of a genuine analytic solution, with an explicit
truncation law near an accumulation exponent?

These are concrete extension questions formulated in this article. They
are not attributed to the repository as verbatim named open problems, and
we do not claim to settle a famous general conjecture in transseries theory.
Theorems~\ref{thm:universal}, \ref{thm:ancestry}, and
\ref{thm:crossover} answer them in the class specified below.

\subsection{What is classical and what is proposed here}
For ordinary integer exponents, reduction of a Fuchsian system to a
resonant polynomial normal form is classical. The entrywise resonance
rule and a product of two constant-matrix powers are recorded, for example,
in Tanny and Yakovenko's discussion of local normal forms
\cite[p.~6]{tanny-yakovenko}. Our construction has that theory as a special
case; neither the Frobenius method nor its integer-exponent normal form is
claimed as a discovery.

General fixed-point methods for generalized series are also established;
van der Hoeven proves an implicit-function theorem for Noetherian
operators \cite{vdh}. We give the particular support and valuation arguments
needed here rather than claiming a new general Hahn fixed-point theorem.
Convergence for Hahn series is itself an established subject, studied in
other settings by Joswig and Smith \cite{joswig-smith}.

The proposed contribution is the combination of an exact semigroup--spectrum
criterion for \emph{universal absolute convergence}, its arbitrarily small
counterexamples, a finite-ancestry refinement of the resonance certificate,
and the explicit crossover analysis of a normalizer at an accumulation
threshold. Full proofs are supplied; historical novelty beyond the sources
inspected remains to be independently assessed. In particular, the article
should not be cited as an independently peer-reviewed breakthrough.

\subsection{Results in one view}
Write $D=x\,d/dx$ and decompose the constant matrix as $A_0=S+N$, where
$S$ is diagonal with real entries and $N$ is nilpotent and commutes with
$S$. The formal construction produces
\begin{equation}\label{eq:intro-solution}
 Y=H(x)x^S\exp(B\log x),\qquad B^m=0.
\end{equation}
The coefficient support of $H-I$ stays in the additive semigroup generated
by that of $\Aplus$. The complete logarithmic filtration is
\begin{equation}\label{eq:intro-filtration}
 \dim_\C\{\text{solutions of logarithmic degree at most }d\}
   =\dim_\C\ker B^{d+1}.
\end{equation}
Moreover, $B$ is a polynomial in finitely many input coefficients once the
constant matrix and exponent data are fixed.

For a fixed positive well-ordered semigroup $\G$, the analytic criterion is
\begin{equation}\label{eq:intro-gap}
 \inf_{\substack{\gamma\in\G_{>0},\ 1\leq i,j\leq m\\
                  \gamma\ne\lambda_i-\lambda_j}}
       \abs{\gamma-(\lambda_i-\lambda_j)}>0.
\end{equation}
There are two essential qualifications. This is a universal assertion over
all coefficient series supported in $\G_{>0}$, not just a specified smaller
generating set. It concerns the normalized gauge constructed here and
absolute generalized-series convergence, not the existence of arbitrary
analytic solutions of the differential equation.

\section{Hahn supports and the differential model}
\label{sec:support}

\subsection{Orientation and allowed supports}
Throughout, $x$ is a positive infinitesimal; analytically it tends to zero
through positive values or lies on a chosen logarithmic cover of a punctured
disk. A Hahn series is written
\[
 f=\sum_{\alpha\in\R}f_\alpha x^\alpha
\]
with well-ordered support in the usual order on $\R$. Thus smaller
exponents are dominant. This is the reverse of the common large-variable
monomial convention, and converting between the two reverses the support
order.

Let $E\subset\R_{>0}$ be nonempty and well ordered, and set
\begin{equation}\label{eq:semigroup}
 \delta=\min E>0,\qquad
 \G=\langle E\rangle
  =\{0\}\cup\bigcup_{k\geq1}
       \{\alpha_1+\cdots+\alpha_k:\alpha_i\in E\}.
\end{equation}
Neither $E$ nor $\G$ is assumed locally finite. For instance,
$E=\{2-1/n:n\geq1\}$ is well ordered but has infinitely many members below
$2$. This type of support is indispensable for the sharpness result.

\begin{lemma}[Support and finite words]\label{lem:words}
The semigroup $\G$ is well ordered. For each $\gamma\in\G$, there are
only finitely many ordered words $(\alpha_1,\ldots,\alpha_k)$ in $E$ whose
sum is $\gamma$. Every such word has length at most
$\lfloor\gamma/\delta\rfloor$.
\end{lemma}
\begin{proof}
The two-set form of Neumann's lemma says that the sum of two well-ordered
subsets of an ordered abelian group is well ordered and that any fixed
element has only finitely many decompositions with one term in each set.
The repository's additive support lemmas provide precisely this two-factor
bridge \cite{repo-wellbased}. Induction therefore proves both assertions
for each fixed word length $k$.

Positivity gives $k\delta\leq\gamma$. Hence only finitely many lengths
can represent a fixed $\gamma$, proving the finite-word assertion.
To prove well ordering of $\G$, let $U\subset\G$ be nonempty and choose
$u\in U$. All elements of $U\cap[0,u]$ belong to the finite union
$\bigcup_{0\leq k\leq\lfloor u/\delta\rfloor}kE$, which is well ordered.
The least element of this intersection is also the least element of $U$.
\end{proof}

The proof uses the Archimedean order of the real exponents in an essential
way: the common positive lower bound controls word length under a finite
cutoff. The analogous assertion with an infinitesimal lower bound in a
higher-rank ordered group requires a different hypothesis.

\subsection{The ambient differential algebra}
Let
\[
 \K=\C((x^\R))
\]
be the Hahn field, and let $\K[L]$ be the polynomial ring in an
algebraically independent variable $L$. Define the strong derivation
\begin{equation}\label{eq:derivation}
 D\!\left(\sum f_\alpha x^\alpha\right)
   =\sum \alpha f_\alpha x^\alpha,\qquad DL=1.
\end{equation}
The multiplication rule is meaningful because the coefficient convolutions
are finite. The Leibniz rule follows by the equality
$\alpha+\beta=\gamma$ in each such convolution. Analytically, $L$ will be
$\log x$.

Let $\mathscr H_\G$ denote the Hahn-series subring with support in $\G$,
and let $\mathfrak m_\G$ be its positive-support ideal. For a nonzero
series or matrix series, $\val(f)$ is the least exponent with nonzero
coefficient, with $\val(0)=+\infty$. Then
\[
 \val(f+g)\geq\min(\val(f),\val(g)),\qquad
 \val(fg)\geq\val(f)+\val(g).
\]
Matrices need the inequality, rather than equality, because leading matrix
products can vanish.

\begin{lemma}[Positive-support limits and inverses]\label{lem:limits}
Suppose $F_n$ are matrix series supported in $\G$ and
$\val(F_{n+1}-F_n)\to+\infty$. Each coefficient eventually stabilizes,
and the coefficientwise limit is a Hahn series supported in $\G$.
If $h\in\Mat_m(\mathfrak m_\G)$, then
\[
 (I+h)^{-1}=\sum_{k\geq0}(-h)^k
\]
is a well-defined matrix Hahn series supported in $\G$.
\end{lemma}
\begin{proof}
Any fixed exponent lies below the valuations of all sufficiently late
differences. The limit support is a subset of the well-ordered set $\G$.
For the inverse, $\val(h^k)\geq k\delta$, so the displayed sum is
coefficientwise finite. Multiplication by $I+h$ telescopes, and the
remaining power tends coefficientwise to zero.
\end{proof}

\begin{lemma}[Constants and the missing Euler residue]\label{lem:constants}
The constant field of $\K$ under $D$ is $\C$, and no derivative in
$\K$ has a nonzero coefficient at $x^0$. The constants of $\K[L]$ are
also exactly $\C$.
\end{lemma}
\begin{proof}
The first two statements follow immediately from
$[x^\alpha]Df=\alpha f_\alpha$. For the third, suppose
$g=\sum_{j=0}^d g_jL^j$ is nonzero and $Dg=0$. Its highest coefficient
satisfies $Dg_d=0$, hence $g_d\in\C$. If $d>0$, the coefficient of
$L^{d-1}$ gives $Dg_{d-1}=-dg_d$, impossible because the right side has a
nonzero $x^0$ coefficient. Thus $d=0$ and $g\in\C$.
\end{proof}

This elementary residue obstruction is the scalar mechanism behind the
logarithmic variable. The matrix certificate below records how different
spectral levels couple that mechanism.

% ed. (2026-09-29): names the repository statements this lemma specializes;
% the article could not read the canonical volume.
\begin{ednote}\label{ed:constants}
Lemma~\ref{lem:constants} is the depth-$0$ case, written for the Euler
derivation $D=x\,d/dx$, of results in the canonical volume
(\path{Analysis/Transseries/docs/series-and-transseries/Transseries_And_Inversion/transseries_and_inversion.tex}),
which this article could not retrieve: since $[x^0](x\,dF/dx)=[x^{-1}]\,dF/dx$,
its residue statement is equation \path{plt:eq:ext-residue} of Theorem
\path{plt:thm:ext-tower-strict} at $r=0$, and its statements on constants
are Corollary \path{plt:cor:dif-kernel}. At depth $n\ge1$ the
residue-obstruction package
\path{Analysis/Transseries/docs/series-and-transseries/Residue_Obstructions_Logarithmic_Depth_Promotion/}
(\texttt{residue\_fredholm.tex}, Theorem \texttt{thm:residue}) sharpens the
same residue to an exact sequence. The volume uses the large-variable
convention; the Euler-form statement is unchanged by $x\mapsto1/x$, which
only changes the sign of $D$.
\end{ednote}

\section{The homological equation and normalized gauge}
\label{sec:normal}

\subsection{The constant matrix and its spectral blocks}
Consider
\begin{equation}\label{eq:system}
 DY=(A_0+\Aplus(x))Y,\qquad
 \Aplus(x)=\sum_{\alpha\in E}A_\alpha x^\alpha.
\end{equation}
Zero coefficients are allowed. Assume $A_0\in\Mat_m(\C)$ has only real
eigenvalues. Choose once and for all a Jordan basis and write
\begin{equation}\label{eq:jordan}
 A_0=S+N,\qquad S=\diag(\lambda_1,\ldots,\lambda_m),\qquad
 [S,N]=0,\quad N^m=0.
\end{equation}
Changing this constant basis changes the normalization convention; all
uniqueness assertions below refer to the chosen basis and projections.

For $\gamma>0$, let $P_\gamma$ be the coordinate projection
\begin{equation}\label{eq:projection}
 (P_\gamma M)_{ij}=
 \begin{cases}M_{ij},&\lambda_i-\lambda_j=\gamma,\\0,&\text{otherwise},\end{cases}
 \qquad Q_\gamma=I-P_\gamma.
\end{equation}
Here the $I$ in $Q_\gamma$ is the identity operator on matrix space.
Define the finite resonance set
\begin{equation}\label{eq:resonances}
 \Res=\{\lambda_i-\lambda_j>0:1\leq i,j\leq m\}\cap\G,
 \qquad r_* = \max(\Res\cup\{0\}).
\end{equation}

The homological operator is
\[
 \mathcal L_\gamma=\gamma I-\ad(A_0),\qquad
 \ad(A_0)M=A_0M-MA_0.
\]
Both $\ad(N)$ and $\mathcal L_\gamma$ preserve every eigenspace of
$\ad(S)$ because $[S,N]=0$.

\begin{lemma}[Explicit nonresonant inverse]\label{lem:inverse}
Define the operator $\Dg_\gamma$ by
\[
 (\Dg_\gamma M)_{ij}=
 \begin{cases}
 \displaystyle\frac{M_{ij}}{\gamma-(\lambda_i-\lambda_j)},
    &\gamma\ne\lambda_i-\lambda_j,\\[6pt]
 0,&\gamma=\lambda_i-\lambda_j.
 \end{cases}
\]
Then
\begin{equation}\label{eq:T}
 T_\gamma=\sum_{k=0}^{2m-2}
       \Dg_\gamma^{\,k+1}(\ad N)^k
\end{equation}
satisfies
\[
 \mathcal L_\gamma T_\gamma=Q_\gamma,
 \qquad P_\gamma T_\gamma=0.
\]
For each $F$, the unique pair $(H,R)$ satisfying
$P_\gamma H=0$, $R=P_\gamma R$, and
$\mathcal L_\gamma H+R=F$ is
$H=T_\gamma F$, $R=P_\gamma F$.
\end{lemma}
\begin{proof}
Left and right multiplication by $N$ commute, so the binomial expansion of
$(\ad N)^{2m-1}$ has in every term either at least $m$ left factors or at
least $m$ right factors. Thus $(\ad N)^{2m-1}=0$. On the range of
$Q_\gamma$, $\gamma I-\ad(S)$ is invertible with inverse
$\Dg_\gamma$, which commutes with $\ad N$. The finite geometric sum
\eqref{eq:T} is therefore the inverse of
$(\gamma I-\ad S)-\ad N$ on that range.

The resonant projection of $\mathcal L_\gamma H$ vanishes when
$P_\gamma H=0$, since the operator preserves the spectral blocks.
Projecting the equation gives $R=P_\gamma F$. The complementary equation
has the stated unique solution.
\end{proof}

The finite geometric sum parallels the repository's explicit inverse of
$\partial_L-c$ on bounded-degree logarithmic polynomials
\cite{repo-block}; here the nilpotent operator is $\ad N$ instead of
polynomial differentiation.

\subsection{The formal normalization theorem}
\begin{theorem}[Normalized Hahn--Fuchsian form]\label{thm:normal}
For \eqref{eq:system}--\eqref{eq:jordan}, there is a unique pair
\begin{equation}\label{eq:HR}
 H=I+\sum_{\gamma\in\G_{>0}}H_\gamma x^\gamma,
 \qquad R(x)=\sum_{\rho\in\Res}R_\rho x^\rho
\end{equation}
with $P_\gamma H_\gamma=0$ and $P_\rho R_\rho=R_\rho$, such that
\begin{equation}\label{eq:gauge}
 DH=(A_0+\Aplus)H-H(A_0+R).
\end{equation}
All series are well-defined Hahn series. No exponent outside $\G$ is
introduced into $H$ or $R$.
\end{theorem}
\begin{proof}
Extend $A_\gamma$ to be zero outside $E$. Equating the coefficient of
$x^\gamma$ in \eqref{eq:gauge} gives
\begin{align}
 F_\gamma={}&A_\gamma+
   \sum_{\substack{\alpha+\beta=\gamma\\\alpha,\beta>0}}
          A_\alpha H_\beta
   -\sum_{\substack{\beta+\eta=\gamma\\\beta,\eta>0}}
          H_\beta R_\eta,                                      \label{eq:F-coeff}\\
 H_\gamma={}&T_\gamma F_\gamma,\qquad
 R_\gamma=P_\gamma F_\gamma.                                  \label{eq:recurrence}
\end{align}
Every coefficient on the right of \eqref{eq:F-coeff}, except the input
$A_\gamma$, has exponent strictly less than $\gamma$. The sums are
finite by Lemma~\ref{lem:words}. Well-ordered recursion and
Lemma~\ref{lem:inverse} give existence and uniqueness. If
$\gamma\notin\Res$, then $P_\gamma=0$, so $R_\gamma=0$. The resulting
supports are subsets of $\G$, proving their admissibility.

For a second construction that will also give an analytic estimate, let
$\Ph$ and $\Th$ act coefficientwise by $P_\gamma$ and $T_\gamma$ on
positive-support matrix series. Set
\begin{equation}\label{eq:fixed-point}
 \Phi(F)=\Aplus+\Aplus(\Th F)-(\Th F)(\Ph F),\qquad F_0=0,
 \quad F_{n+1}=\Phi(F_n).
\end{equation}
All $F_n$ have support in $\G_{>0}$. Both coefficientwise operators
preserve or increase valuation. If $F,G$ have positive support, then
\[
 \Phi(F)-\Phi(G)=\Aplus\Th(F-G)
  -\Th(F-G)\Ph F-(\Th G)\Ph(F-G).
\]
Consequently,
\begin{equation}\label{eq:valuation-contraction}
 \val(\Phi(F)-\Phi(G))\geq\val(F-G)+\delta.
\end{equation}
It follows by induction that
$\val(F_{n+1}-F_n)\geq(n+1)\delta$. Lemma~\ref{lem:limits} gives a
limit $F$; finite coefficient convolutions show $F=\Phi(F)$. Take
$H=I+\Th F$ and $R=\Ph F$. Lemma~\ref{lem:inverse} yields
\eqref{eq:gauge}. A difference of two fixed points would have its finite
valuation increased by at least $\delta$, a contradiction. This also
proves uniqueness without invoking transfinite recursion separately.
\end{proof}

\begin{corollary}[Finite iteration count for the normal form]
\label{cor:iterations}
The iterates in \eqref{eq:fixed-point} satisfy
\[
 \val(F-F_n)\geq(n+1)\delta.
\]
In particular, $R=\Ph F_n$ as soon as $(n+1)\delta>r_*$.
\end{corollary}
\begin{proof}
All differences in the tail beginning with $F_{n+1}-F_n$ have valuation
at least $(n+1)\delta$. The projection $\Ph$ retains only exponents at
most $r_*$.
\end{proof}

\begin{remark}
A finite number of formal iterations does not mean that each whole
iterate has finite support. When exponents accumulate, even its part below
$r_*$ can be infinite. The genuinely finite input certificate is proved
in Section~\ref{sec:ancestry}, not by silently replacing well ordering
with local finiteness.
\end{remark}

\section{A finite matrix controls every logarithm}
\label{sec:logs}

\begin{theorem}[Fundamental matrix and logarithmic filtration]
\label{thm:logs}
For the normalized pair $(H,R)$, put
\begin{equation}\label{eq:B}
 B=N+\sum_{\rho\in\Res}R_\rho.
\end{equation}
Then $B^m=0$, and
\begin{equation}\label{eq:Y}
 Y=H(x)x^S e^{BL},\qquad
 x^S=\diag(x^{\lambda_1},\ldots,x^{\lambda_m}),\quad
 e^{BL}=\sum_{k=0}^{m-1}\frac{B^kL^k}{k!}
\end{equation}
is a fundamental matrix over $\K[L]$. Every solution vector with entries
in $\K[L]$ is uniquely $Yc$ for $c\in\C^m$.

For each integer $d\geq0$, the space of solutions of logarithmic degree
at most $d$ has dimension
\begin{equation}\label{eq:filtration}
 \dim_\C\Sol_{\leq d}=\dim_\C\ker B^{d+1}.
\end{equation}
A logarithm-free fundamental matrix exists in this class if and only if
$B=0$. If $\nu(B)$ is the least positive integer with $B^{\nu(B)}=0$
(with $\nu(0)=1$), then every fundamental matrix in $\K[L]$ has maximum
logarithmic degree exactly $\nu(B)-1\leq m-1$.
\end{theorem}
\begin{proof}
The matrix $N$ preserves each eigenspace of $S$, whereas $R_\rho$ maps
the $\lambda$ eigenspace into the $\lambda+\rho$ eigenspace. Order the
distinct eigenvalues increasingly. The matrix $B$ is block triangular,
with nilpotent diagonal blocks inherited from $N$. Its characteristic
polynomial is $z^m$, so the Cayley--Hamilton theorem gives $B^m=0$.

The spectral support rule gives
\begin{equation}\label{eq:conjugateB}
 x^S Bx^{-S}=N+\sum_{\rho\in\Res}x^\rho R_\rho=N+R(x).
\end{equation}
Differentiating $x^S e^{BL}$, using $DL=1$, yields
\[
 D(x^S e^{BL})=(S+N+R)x^S e^{BL}=(A_0+R)x^S e^{BL}.
\]
Combining this with \eqref{eq:gauge} proves that $Y$ solves
\eqref{eq:system}. The inverse
\[
 Y^{-1}=e^{-BL}x^{-S}H^{-1}
\]
has entries in $\K[L]$ by Lemma~\ref{lem:limits}; hence $Y$ is a
fundamental matrix.

If $y$ is any other solution in this algebra, then
$D(Y^{-1}y)=0$. Lemma~\ref{lem:constants} gives $Y^{-1}y=c\in\C^m$.
The coefficient of $L^k$ in $Yc$ is $Hx^SB^kc/k!$. Because $Hx^S$ is
invertible and independent of $L$, this coefficient is zero exactly when
$B^kc=0$. Thus $\deg_L(Yc)\leq d$ is equivalent to $B^{d+1}c=0$,
which proves \eqref{eq:filtration}.

Any fundamental matrix is $YC$ with $C\in\mathrm{GL}_m(\C)$. Its top
nonzero logarithmic coefficient is
$Hx^SB^{\nu(B)-1}C/(\nu(B)-1)!$, which cannot vanish. The remaining
assertions follow.
\end{proof}

\begin{corollary}[Invariance under log-free gauges]
Multiplication of a solution vector by an invertible matrix over $\K$
preserves its logarithmic degree. Hence the dimensions in
\eqref{eq:filtration} are invariants of log-free gauge equivalence, even
though the entries of the normalized matrix $B$ depend on the chosen
normalization.
\end{corollary}
\begin{proof}
An invertible log-free matrix cannot kill the highest nonzero coefficient
of a polynomial in $L$.
\end{proof}

\begin{example}[The bound $m-1$ is attained]\label{ex:chain}
Let $S=\diag(0,1,\ldots,m-1)$, $N=0$, and
$J=\sum_{j=1}^{m-1}E_{j+1,j}$, where $E_{ij}$ is a matrix unit.
For $\Aplus=xJ$, all entries are resonant. The normalized gauge is
$H=I$, the normal form is $R=xJ$, and $B=J$. Since
$J^{m-1}\ne0$ and $J^m=0$, logarithmic degree $m-1$ is unavoidable.
The result is an exact identity, not a merely asymptotic model.
\end{example}

\begin{remark}[The ambient algebra matters]
The assertion is about polynomial dependence on the one indeterminate
$L$ over the log-free Hahn field $\K$. It does not identify $L$ with an
element already in $\K$, nor does it claim that every possible analytic
representation must visibly display a logarithm. An unrestricted special
function can hide one. The filtration is meaningful because the coefficient
algebra and allowed log-free gauges have been fixed.
\end{remark}

\section{Finite ancestry despite infinitely many low exponents}
\label{sec:ancestry}

Define the critical ancestry set
\begin{equation}\label{eq:critical-set}
 E_{\mathrm{crit}}=
 \bigcup_{\rho\in\Res}
 \bigcup_{\substack{k\geq1,\ \alpha_1,\ldots,\alpha_k\in E\\
                       \alpha_1+\cdots+\alpha_k=\rho}}
       \{\alpha_1,\ldots,\alpha_k\}.
\end{equation}
The union over an empty resonance set is empty.

\begin{theorem}[Finite-ancestry certificate]\label{thm:ancestry}
The set $E_{\mathrm{crit}}$ is finite. For fixed $S,N,E$, each matrix
$R_\rho$ is an ordinary polynomial in the entries of
$\{A_\alpha:\alpha\in E_{\mathrm{crit}}\}$, of degree at most
$\lfloor\rho/\delta\rfloor$. Its scalar coefficients depend on the
constant matrix and the finitely many intermediate exponent differences.
In particular, all logarithmic obstruction data are independent of input
coefficients outside $E_{\mathrm{crit}}$.

For a family obtained by varying these finitely many matrix entries,
\begin{equation}\label{eq:algebraic-locus}
 \{\text{maximum logarithmic degree}\leq d\}
       =\{B^{d+1}=0\}
\end{equation}
is an algebraic locus. The ranks of the powers of $B$, and hence the entire
logarithmic filtration, are determined by polynomial minors.
\end{theorem}
\begin{proof}
There are finitely many resonances. Lemma~\ref{lem:words} gives finitely
many words ending at each one, so \eqref{eq:critical-set} is finite.

Expand the recurrence \eqref{eq:F-coeff}--\eqref{eq:recurrence} by
repeated substitution. The operators $T_\gamma$ and $P_\gamma$ are
linear and do not shift exponents. Every product term therefore has input
leaves $A_{\alpha_1},\ldots,A_{\alpha_k}$ whose exponent sum is the
output exponent. Each further substitution replaces a positive exponent
by positive smaller exponents. The expansion ends after at most
$\lfloor\gamma/\delta\rfloor$ input leaves. Finiteness of the possible
words also gives finiteness of all partial exponent sums, so the expression
at each fixed output exponent is an ordinary finite polynomial, not an
infinite coefficient formula.

At output $\rho$, all leaves belong to $E_{\mathrm{crit}}$ by
definition. The only divisions are the nonzero scalar denominators in
$T_\gamma$, which depend on $S$ and intermediate exponents rather than on
input coefficient variables. Thus the polynomial degree is bounded by the
number of leaves. Since $B=N+\sum R_\rho$, its entries are polynomials in
finite data. Theorem~\ref{thm:logs} proves \eqref{eq:algebraic-locus};
the standard minor criteria for matrix rank prove the last assertion.
\end{proof}

\subsection{What the finite statement does not assert}
For an arbitrary infinitely presented set of real exponents, it can be
impossible to decide whether a proposed number is a finite sum of members.
The theorem supplies finite mathematical dependence; it does not supply an
unconditional terminating algorithm for discovering the critical set from
an arbitrary oracle. The accompanying implementation restricts to finite
rational input supports, where exact equality and bounded enumeration are
decidable.

There is a useful distinction between three cutoffs. The largest resonance
$r_*$ bounds where normal-form output can occur. The number
$\lfloor r_*/\delta\rfloor$ bounds the number of input factors in an
obstruction. Neither statement says that all exponents below $r_*$ must be
enumerated. The critical ancestry set selects just the finitely many
exponents that can actually reach a resonance by addition.

\subsection{Relation to the earlier residue-matrix method}
The repository's residue-obstruction article already uses a positive
exponent decrement to bound perturbative paths between finitely many real
roots \cite{repo-residue}. Our finite-ancestry theorem is aligned with that
principle. The distinction is the present matrix normal form, the complete
log-degree filtration via powers of $B$, and especially the analytic
criterion in the next section. The previous scalar fixed-depth problem and
the present universal absolute-convergence problem are not interchangeable.

% ed. (2026-09-29): records which questions of the residue-obstruction
% article this article answers without citing them, and the checked
% consistency between the two.
\begin{ednote}\label{ed:residue-questions}
This article knew the residue-obstruction package only through its README.
Its source,
\path{Analysis/Transseries/docs/series-and-transseries/Residue_Obstructions_Logarithmic_Depth_Promotion/residue_fredholm.tex},
poses in its Section~11 two questions that Theorem~\ref{thm:logs} answers
for the present class. ``The exact minimal power of the promoted
logarithm'' asks for residue matrices whose ranks recover the full
filtration by degree in the new logarithm: here, for homogeneous
solutions, it is $\dim\Sol_{\le d}=\dim\ker B^{d+1}$ (particular solutions
of forced equations are not treated). ``Matrix systems and nonreal
indicial roots'' asks which invariant controls the new-logarithm degree:
here it is the nilpotency index $\nu(B)$, which combines Jordan chains in
$N$ with passages between root levels through the $R_\rho$ (real spectrum
only). The two articles are consistent but not nested: that article works
at depth $n\ge1$ with scalar operators $P(xD)+\sum a_j(xD)^j$, this one at
depth $0$ with first-order systems. At $d=0$, $\dim\ker B=m-\rank B$ is the
analogue of its kernel dimension $r-\rank M$ (Theorem \texttt{thm:fredholm});
Example~\ref{ex:chain} is the analogue of its Jordan-shift family (Theorem
\texttt{thm:sharp-matrix}); and when $N=0$, $B$ is strictly block-triangular
over the $s$ distinct eigenvalues, so the logarithmic degree is at most
$s-1$, its homogeneous bound (Theorem \texttt{thm:promotion}). With
$N\ne0$ the degree can reach $m-1>s-1$, which is no contradiction: there
the multiplicity logarithms already lie in the base field.
\end{ednote}

\section{A sharp criterion for universal analytic normalization}
\label{sec:analytic}

\subsection{The absolute generalized-series algebra}
Use the induced column-sum matrix norm
\[
 \norm{M}_1=\max_j\sum_i|M_{ij}|,
\]
so $\norm I_1=1$, matrix multiplication is submultiplicative, and every
coordinate mask $P_\gamma$ or $Q_\gamma$ is contractive. For $r>0$ set
\begin{equation}\label{eq:analytic-norm}
 \norm F_r=\sum_{\gamma\in\G_{>0}}\norm{F_\gamma}_1r^\gamma.
\end{equation}
A series is called absolutely convergent at radius $r$ when this norm is
finite. Together with a constant term, these series form a Banach algebra
under convolution. A well-ordered subset of $\R$ is countable, so the
notation is an ordinary countable absolute sum; equivalently it can be
read as a supremum over finite subsums.

For real exponents and a chosen branch, $|x^\gamma|=|x|^\gamma$.
Absolute convergence therefore gives holomorphic evaluation on the
logarithmic cover of $0<|x|<r$. The Euler derivative converges normally on
every smaller radius $r'<r$, because
$\sup_{\gamma\geq0}\gamma(r'/r)^\gamma<\infty$. Products and the
Euler derivative may then be evaluated termwise.

Define
\begin{equation}\label{eq:eta}
 \eta(\G,S)=
 \inf\left\{
 |\gamma-(\lambda_i-\lambda_j)|:
 \begin{array}{l}
 \gamma\in\G_{>0},\ 1\leq i,j\leq m,\\[-2pt]
 \gamma\ne\lambda_i-\lambda_j
 \end{array}\right\}.
\end{equation}
The exclusion of exact resonances is crucial: exact resonances are put in
$R$, not divided by.

\begin{lemma}[Which accumulations matter]\label{lem:gap}
The number $\eta(\G,S)$ is positive if and only if no positive
eigenvalue difference is a limit of a strictly increasing sequence of
nonresonant exponents in $\G$.
\end{lemma}
\begin{proof}
Such an increasing sequence forces the infimum to be zero. Conversely,
if the infimum is zero, finiteness of the set of eigenvalue differences
provides one difference $\rho$ and distinct exponents
$\gamma_n\in\G_{>0}$ tending to $\rho$, none equal to it. Since
$\gamma_n\geq\delta$, the difference $\rho$ is positive. A subsequence
approaching $\rho$ from above would contain an infinite strictly decreasing
subsequence, contrary to well ordering of $\G$. There are therefore
infinitely many terms below $\rho$, from which one extracts a strictly
increasing subsequence tending to $\rho$.
\end{proof}

\begin{theorem}[Universal convergence: necessary and sufficient]
\label{thm:universal}
Fix a nontrivial semigroup $\G\subset\R_{\geq0}$ generated as in
\eqref{eq:semigroup}, and fix $A_0=S+N$ as in \eqref{eq:jordan}. The
following statements are equivalent.
\begin{enumerate}[label=\textup{(\roman*)}]
\item $\eta(\G,S)>0$.
\item For every matrix series $\Aplus$ supported in $\G_{>0}$ that is
absolutely convergent at some radius, the normalized formal gauge $H$ of
Theorem~\ref{thm:normal} is absolutely convergent at some positive radius.
\end{enumerate}
If \textup{(i)} fails, counterexamples in \textup{(ii)} can have arbitrarily
small norm at radius $1$ and can be chosen proportional to one square-zero
rank-one constant matrix. Their coefficient series are absolutely
convergent at every positive radius.
\end{theorem}

\subsection{Proof of sufficiency with quantitative bounds}
Assume $\eta>0$ and put
\begin{equation}\label{eq:C}
 C=\sum_{k=0}^{2m-2}\frac{(2\norm N_1)^k}{\eta^{k+1}}.
\end{equation}
The powers are ordinary nonnegative real numbers; for $N=0$, the formula
reduces to $C=1/\eta$. Since
$\norm{\ad(N)M}_1\leq2\norm N_1\norm M_1$, the finite inverse
\eqref{eq:T} gives
\begin{equation}\label{eq:operator-bounds}
 \norm{\Th F}_r\leq C\norm F_r,\qquad
 \norm{\Ph F}_r\leq\norm F_r.
\end{equation}
Suppose $a_0=\norm{\Aplus}_{r_0}<\infty$. For $0<r\leq r_0$,
\begin{equation}\label{eq:radius-decay}
 a:=\norm{\Aplus}_r\leq a_0(r/r_0)^\delta.
\end{equation}
Choose $r>0$ with $Ca\leq1/8$. If $a_0>0$, one explicit permissible
choice is
\begin{equation}\label{eq:explicit-radius}
 r=r_0\min\left\{1,(8Ca_0)^{-1/\delta}\right\}.
\end{equation}
The zero-input case is immediate.

On the closed ball $\norm F_r\leq2a$, the same map $\Phi$ as in
\eqref{eq:fixed-point} satisfies
\begin{align*}
 \norm{\Phi(F)}_r
 &\leq a+aC(2a)+C(2a)(2a)
 =a+6Ca^2\leq\tfrac74a<2a,\\
 \norm{\Phi(F)-\Phi(G)}_r
 &\leq(Ca+4Ca)\norm{F-G}_r
 =5Ca\norm{F-G}_r.
\end{align*}
Thus $q:=5Ca\leq5/8$ is a contraction constant. The Banach fixed-point
theorem yields an absolutely convergent $F$ in this ball. Its coefficient
identity is the formal fixed-point equation, so formal uniqueness identifies
it with the Hahn solution constructed earlier. Consequently,
\begin{equation}\label{eq:analytic-estimates}
 \norm{H-I}_r\leq2Ca\leq\tfrac14,
 \qquad\norm R_r\leq2a,
 \qquad\norm{H^{-1}}_r\leq\frac1{1-2Ca},
\end{equation}
where the last norm includes the constant term.

The Picard iterates provide the additional error certificate
\begin{equation}\label{eq:picard-bound}
 \norm{F-F_n}_r\leq\frac{a q^n}{1-q},\qquad
 \norm{H-(I+\Th F_n)}_r\leq\frac{Ca q^n}{1-q}.
\end{equation}
Indeed, the first iterate is $F_1=\Aplus$ and the differences decrease by
at least the factor $q$. This proves the analytic assertion and supplies a
constructive bound rather than just a qualitative convergence theorem.

All identities now hold analytically on smaller punctured radii. In
particular, \eqref{eq:Y} is an actual fundamental matrix on any fixed
logarithmic branch there. The word ``Fuchsian'' here refers to the Euler
form and limiting constant coefficient; we do not assume that
$\Aplus$ has an ordinary Taylor expansion at zero.

\subsection{Proof of necessity and small counterexamples}
Assume $\eta=0$. By Lemma~\ref{lem:gap}, choose a positive difference
$\rho=\lambda_i-\lambda_j$ and a strictly increasing sequence
$\gamma_n\in\G_{>0}$ with limit $\rho$, with no term equal to
$\rho$. Passing to a subsequence, arrange
\begin{equation}\label{eq:dn}
 d_n:=\rho-\gamma_n\leq2^{-n}.
\end{equation}
Choose a nonzero right eigenvector $u$ of $A_0$ at $\lambda_i$ and a
nonzero left eigenvector $v^T$ at $\lambda_j$. Normalize
$U=uv^T$ so that $\norm U_1=1$. Distinct eigenvalues imply $v^Tu=0$,
so
\begin{equation}\label{eq:rank-one}
 U^2=0,\qquad [A_0,U]=[S,U]=\rho U,\qquad[N,U]=0.
\end{equation}
For any $\varepsilon>0$ define
\begin{equation}\label{eq:counter-general}
 \Aplus(x)=\varepsilon\sum_{n\geq1}\frac{d_n}{n}x^{\gamma_n}U.
\end{equation}
The exponents remain in a bounded interval, and
$\sum d_n/n\leq\sum2^{-n}/n<\infty$. Thus the input is absolutely
convergent at every radius, with norm at radius $1$ tending to zero as
$\varepsilon\to0$.

The normalized formal gauge is exactly
\begin{equation}\label{eq:counter-H}
 H=I-\varepsilon\sum_{n\geq1}\frac{x^{\gamma_n}}nU,
 \qquad R=0.
\end{equation}
To check this, write $H=I+hU$. The nonlinear products vanish because
$U^2=0$, and the gauge equation reduces to
$(D-\rho)h=\varepsilon\sum d_nx^{\gamma_n}/n$. Its coefficient
solution is $h_{\gamma_n}=-\varepsilon/n$. All these coefficients are
nonresonant, so the gauge is normalized; uniqueness proves the assertion.
For every $r>0$, the factors $r^{\gamma_n}$ tend to $r^\rho>0$.
Consequently
$\sum r^{\gamma_n}/n=\infty$, and the normalized gauge is not
absolutely convergent at any radius. This proves necessity and the
smallness statement in Theorem~\ref{thm:universal}.\hfill$\square$

\subsection{Consequences and exact boundaries}
\begin{corollary}[Finite positive generators need no arithmetic condition]
\label{cor:finite-grid}
If $\G$ is generated by finitely many positive real numbers, every
absolutely convergent $\G$-supported coefficient series has a convergent
normalized gauge. Rational relations among the generators are irrelevant.
\end{corollary}
\begin{proof}
Below any fixed bound, the nonnegative integer coefficient of each positive
generator is bounded. Hence $\G$ is locally finite. No finite positive
eigenvalue difference can be an accumulation point.
\end{proof}

For example, $\langle1,\sqrt2\rangle$ is locally finite even though the
additive \emph{group} generated by $1$ and $\sqrt2$ is dense in $\R$.
Replacing the positive semigroup by that group would give the wrong
small-divisor diagnosis. The dangerous event is not irrationality, nor
non-local-finiteness by itself: it is a semigroup accumulation specifically
at a positive spectral difference.

In dimension one, there is no positive eigenvalue difference and
$\eta\geq\delta$. Thus rank two is the smallest dimension for the
failure mechanism. Exact resonances and near-resonances also have opposite
roles. Exact resonances produce finitely many logarithmic data. An infinite
sequence of near-resonances can destroy convergence while producing no
additional exact resonant term at all.

\begin{remark}[Universal quantifiers cannot be weakened silently]
The counterexample uses a sequence from $\G$, which need not lie in a
prescribed proper generating subset $E$. The theorem quantifies over all
inputs supported in $\G_{>0}$. For one fixed input, $\eta=0$ does not
force divergence: its coefficients may decay enough, or matrix products
may suppress the bad divisors. A theorem quantifying only over series
supported in a smaller $E$ is a separate problem.
\end{remark}

\section{Indirect resonance and a finite-data example}
\label{sec:indirect}

A direct inspection of the original coefficient at a resonant exponent is
insufficient: earlier nonresonant terms can generate a logarithm. This is
already visible in a two-dimensional system with ordinary integer powers.

Let
\begin{equation}\label{eq:indirect-input}
 S=\begin{pmatrix}0&0\\0&2\end{pmatrix},\quad N=0,\quad
 A_1=\begin{pmatrix}a&b\\c&d\end{pmatrix},\quad
 A_2=\begin{pmatrix}0&0\\e&0\end{pmatrix}.
\end{equation}
The only positive resonance is $2$, in the lower-left entry. At exponent
$1$, the recurrence gives
\[
 H_1=\begin{pmatrix}a&b/3\\-c&d\end{pmatrix},\qquad R_1=0.
\]
At exponent $2$, the lower-left entry of $A_2+A_1H_1$ is
$e+c(a-d)$. Therefore
\begin{equation}\label{eq:indirect-B}
 B=\bigl(e+c(a-d)\bigr)E_{21}.
\end{equation}
The exact no-log condition is the algebraic equation
$e+c(a-d)=0$, not merely $e=0$. When it fails, the log-free solution
subspace has dimension one, and a fundamental matrix has logarithmic
degree one.

For a fully explicit example, take $a=c=1$ and $b=d=e=0$. The system is
\begin{equation}\label{eq:indirect-concrete}
 Dy_1=xy_1,\qquad Dy_2=2y_2+xy_1.
\end{equation}
A fundamental matrix is
\begin{equation}\label{eq:indirect-Y}
 Y(x)=
 \begin{pmatrix}
 e^x&0\\[2pt]
 \displaystyle -x+x^2\log x+
    \sum_{k=2}^{\infty}\frac{x^{k+1}}{(k-1)k!}&x^2
 \end{pmatrix}.
\end{equation}
The series is entire. Direct application of $D-2$ to its lower-left entry
gives $xe^x$, including the $x^2$ term supplied by $x^2\log x$.
The normalized gauge has $H_{22}=1$, $H_{12}=0$ and
\[
 H_{11}=e^x,\qquad
 H_{21}=-x+\sum_{n=3}^{\infty}\frac{x^n}{(n-2)(n-1)!}.
\]
In particular its resonant $x^2E_{21}$ coefficient vanishes, as required,
and $B=E_{21}$.

Now allow an arbitrary extra tail with exponents in
\[
 \{3-1/n:n\geq2\}.
\]
Every new exponent exceeds $2$. Its coefficient matrices may be arbitrary
Hahn data: none can appear in a word ending at the sole resonance $2$.
Thus the obstruction \eqref{eq:indirect-B} is unchanged. The augmented
semigroup has an accumulation at $3$, but this is not a positive spectral
difference for this system. There are only finitely many semigroup elements
in a sufficiently small neighborhood of $2$, so the uniform divisor gap
remains positive. Absolutely convergent tails therefore also preserve
analytic normalizability. This example separates support accumulation,
finite logarithmic data, and analytic failure in a single calculation.

\section{An explicit summability threshold at an accumulation exponent}
\label{sec:model}

\subsection{The family and its divergent formal normalizer}
Let $p>1$ and
\begin{equation}\label{eq:ap}
 a_p(x)=\sum_{n=1}^{\infty}\frac{x^{2-1/n}}{n^p},\qquad
 S=\begin{pmatrix}2&0\\0&0\end{pmatrix},\qquad U=E_{12}.
\end{equation}
Consider
\begin{equation}\label{eq:p-system}
 DY=(S+a_p(x)U)Y.
\end{equation}
The input series is absolutely convergent at every positive radius:
its exponents belong to $[1,2)$ and $\sum n^{-p}$ converges. It is
holomorphic on the logarithmic cover and tends to zero as $x\to0^+$.
Its well-ordered support approaches the positive spectral difference $2$
from below.

\begin{proposition}[Exact threshold]\label{prop:threshold}
The normalized formal gauge and resonant matrix are
\begin{equation}\label{eq:hhatp}
 \widehat H_p=I+\widehat h_p U,
 \qquad\widehat h_p=-\sum_{n=1}^{\infty}n^{1-p}x^{2-1/n},
 \qquad R=0,
\end{equation}
and $B=0$. The gauge is absolutely convergent at some positive radius if
and only if $p>2$; in that case it converges at every positive radius.
\end{proposition}
\begin{proof}
Since $[S,U]=2U$ and $U^2=0$, the coefficient equation is
$(D-2)\widehat h_p=a_p$. Dividing the coefficient at $2-1/n$ by
$-1/n$ gives $-n^{1-p}$. There are no product terms and no input at the
resonance itself, so $R=0$. At any fixed positive radius, $r^{2-1/n}$
tends to $r^2$, so absolute convergence is equivalent to that of
$\sum n^{1-p}$.
\end{proof}

Although the generated semigroup contains $2=1+1$, the exact resonant
coefficient remains zero: the only possible matrix product is a multiple
of $U^2=0$. This illustrates why membership of a resonance in the
semigroup alone is not the obstruction certificate.

\subsection{Renormalized analytic solutions for every $p>1$}
The formal series can diverge while the actual differential equation still
has a particularly explicit solution. For $t=-\log x$ define
\begin{align}
 \Psi_p(t)
   &=\sum_{n=1}^{\infty}n^{1-p}(e^{t/n}-1),                       \label{eq:Psi}\\
 h_p(x)&=-x^2\Psi_p(-\log x).                                   \label{eq:hp}
\end{align}
The subtraction is made \emph{before} summation.

\begin{theorem}[Analytic realization by renormalization]
\label{thm:renorm}
For every $p>1$, $\Psi_p$ is entire and
\begin{equation}\label{eq:zeta-expansion}
 \Psi_p(t)=\sum_{k=1}^{\infty}\zeta(p+k-1)\frac{t^k}{k!}.
\end{equation}
The function $h_p$ is holomorphic on the logarithmic cover, satisfies
$(D-2)h_p=a_p$ and $h_p(1)=0$, and gives the fundamental matrix
\begin{equation}\label{eq:actual-p}
 Y_p=(I+h_pU)x^S.
\end{equation}
For $p>2$, it agrees with the convergent formal normalizer up to the
homogeneous term
\begin{equation}\label{eq:p-comparison}
 h_p=\widehat h_p+\zeta(p-1)x^2.
\end{equation}
For $1<p\leq2$, the defining difference series still converges, but its
two pieces cannot be separately summed absolutely.
\end{theorem}
\begin{proof}
For $|t|\leq T$,
\[
 n^{1-p}|e^{t/n}-1|\leq T e^T n^{-p}.
\]
The series \eqref{eq:Psi} therefore converges uniformly on compact sets,
proving that $\Psi_p$ is entire. Expanding each exponential and using
absolute convergence gives
\[
 \sum_{n\geq1}\sum_{k\geq1}
   \frac{|t|^k}{k!n^{p+k-1}}
 \leq\zeta(p)(e^{|t|}-1)<\infty.
\]
Interchanging the sums proves \eqref{eq:zeta-expansion}.

Since $D=-d/dt$ on functions of $t$ and $x^2=e^{-2t}$,
\[
 (D-2)(-x^2\Psi_p(t))=x^2\Psi_p'(t)
   =x^2\sum_{n\geq1}n^{-p}e^{t/n}=a_p(x).
\]
The value at $x=1$ is zero. This scalar equation and $U^2=0$ verify
\eqref{eq:actual-p} directly. If $p>2$, one can split
\eqref{eq:Psi} into two absolutely convergent sums; their difference
is precisely \eqref{eq:p-comparison}. If $p\leq2$, the series of
subtracted constants is $\sum n^{1-p}$ and diverges, so that split is
not legitimate.
\end{proof}

\subsection{Every finite asymptotic prefix is nevertheless correct}
For $N\geq1$, set
\begin{equation}\label{eq:hN}
 h_{p,N}(x)=-\sum_{n=1}^{N}n^{1-p}(x^{2-1/n}-x^2),\qquad
 S_{p,N}(t)=\sum_{n>N}n^{1-p}(e^{t/n}-1).
\end{equation}
For $0<x<1$, $t=-\log x>0$ and
\begin{equation}\label{eq:error-exact}
 h_p-h_{p,N}=-x^2 S_{p,N}(t).
\end{equation}

\begin{lemma}[Elementary two-sided tail certificate]\label{lem:tail}
For $p>1$, $N\geq1$, $t>0$,
\begin{equation}\label{eq:tail-bounds}
 \frac{t}{p-1}(N+1)^{1-p}
 \leq S_{p,N}(t)
 \leq\frac{t e^{t/(N+1)}}{p-1}N^{1-p}.
\end{equation}
\end{lemma}
\begin{proof}
For $u\geq0$, $u\leq e^u-1\leq ue^u$. Thus
\[
 t\sum_{n>N}n^{-p}\leq S_{p,N}(t)
   \leq t e^{t/(N+1)}\sum_{n>N}n^{-p}.
\]
The integral test bounds the final sum between
$(N+1)^{1-p}/(p-1)$ and $N^{1-p}/(p-1)$.
\end{proof}

\begin{proposition}[Asymptotic realization without absolute summation]
\label{prop:prefix}
For every fixed $N\geq1$ and every $p>1$,
\begin{equation}\label{eq:prefix}
 h_p(x)+\sum_{n=1}^{N}n^{1-p}x^{2-1/n}
 \sim -(N+1)^{1-p}x^{2-1/(N+1)}
 \quad(x\to0^+).
\end{equation}
In particular, $h_p$ realizes all finite Poincare prefixes of
$\widehat h_p$, including in the divergent range $1<p\leq2$.
Adding any constant multiple of $x^2$ leaves every such prefix unchanged.
\end{proposition}
\begin{proof}
For fixed $N$, isolate the first term of the tail:
\[
 S_{p,N}(t)=(N+1)^{1-p}(e^{t/(N+1)}-1)+S_{p,N+1}(t).
\]
The upper estimate in Lemma~\ref{lem:tail} shows that the last term is
$O(t e^{t/(N+2)})$, which is little-oh of $e^{t/(N+1)}$.
Consequently $S_{p,N}(t)\sim(N+1)^{1-p}e^{t/(N+1)}$.
The left side of \eqref{eq:prefix} equals
$x^2\sum_{n\leq N}n^{1-p}-x^2S_{p,N}(t)$. The finite constant sum is
negligible compared with $e^{t/(N+1)}$, proving the formula.
Finally, $x^2=o(x^{2-1/n})$ for every fixed $n$.
\end{proof}

This is a precise boundary between three assertions often conflated in
transseries manipulations. The series is formally valid. It is asymptotically
realized to each fixed finite order. It is absolutely summable only for
$p>2$. Neither of the first two assertions implies the third. The
undetectable $Cx^2$ term is also the natural analogue, for this accumulating
scale, of the repository's flat-function ambiguity.

\section{The truncation crossover and an explicit uniform remainder}
\label{sec:crossover}

Fixed $N$ describes successive asymptotic prefixes; $N\to\infty$ at
fixed $x$ describes convergence of the renormalized sum. The interesting
transition occurs when both grow and $N$ is proportional to
$t=\log(1/x)$. Set
\begin{equation}\label{eq:fPhi}
 f_p(s)=s^{1-p}(e^{1/s}-1),\qquad
 \Phi_p(c)=\int_c^\infty f_p(s)\,ds
          =\int_0^{1/c}u^{p-3}(e^u-1)\,du.
\end{equation}
Both integrals converge for $p>1$ and $c>0$.

\begin{theorem}[Certified accumulation-scale crossover]
\label{thm:crossover}
Let $p>1$, $t>0$, $N\geq1$, and $c=N/t$. Then
\begin{equation}\label{eq:crossover}
 S_{p,N}(t)=t^{2-p}
 \left(\Phi_p(c)-\frac{f_p(c)}{2t}
             -\frac{f_p'(c)}{12t^2}+\mathcal E_{p,N,t}\right),
\end{equation}
with the explicit bound
\begin{equation}\label{eq:EM-bound}
 |\mathcal E_{p,N,t}|
       \leq\frac{|f_p'''(c)|}{360t^4}.
\end{equation}
In particular, the expansion is uniform when $c$ ranges over a compact
subinterval of $(0,\infty)$, with $p$ fixed. If $c>0$ is fixed and
$t\to\infty$ along values with $ct\in\N$, then
\begin{equation}\label{eq:EM-limit}
 t^4\mathcal E_{p,ct,t}\longrightarrow\frac{f_p'''(c)}{720}.
\end{equation}
\end{theorem}
\begin{proof}
From the definition,
\[
 S_{p,N}(t)=t^{1-p}\sum_{n>N} f_p(n/t).
\]
Thus $t^{p-2}S_{p,N}(t)$ is a right-endpoint Riemann sum with mesh
$h=1/t$ beginning strictly after $c=Nh$.

We record the needed half-line Euler--Maclaurin formula, including the
constants. For a smooth integrable function $f$ whose relevant derivatives are
integrable and decay at infinity,
\begin{align}
 h\sum_{k=1}^{\infty}f(c+kh)
   ={}&\int_c^\infty f(s)\,ds-\frac h2f(c)
      -\frac{h^2}{12}f'(c)+\frac{h^4}{720}f'''(c)
      +\mathcal R_4,                                         \label{eq:EM}\\
 |\mathcal R_4|\leq{}&\frac{h^4}{720}
                          \int_c^\infty|f^{(4)}(s)|\,ds.     \label{eq:EM-rem}
\end{align}
This follows by integrating by parts with the periodic Bernoulli
polynomials on successive mesh intervals. The coefficients are
$B_2/2!=1/12$ and $B_4/4!=-1/720$; the boundary terms at infinity vanish.
The periodic fourth Bernoulli polynomial divided by $4!$ has absolute
value at most $1/720$, giving \eqref{eq:EM-rem}. This also proves the
version on an infinite interval by first using a finite number of mesh
intervals and passing to the limit.

For our function,
\begin{equation}\label{eq:complete-monotone}
 f_p(s)=\sum_{k=1}^{\infty}\frac{s^{-(p+k-1)}}{k!}.
\end{equation}
The series and each derivative converge uniformly on $[c,\infty)$ after
the appropriate integrable power bound. Every summand is completely
monotone, so $f_p^{(4)}\geq0$ and $f_p'''\leq0$. Its derivatives tend
to zero at infinity; hence
\[
 \int_c^\infty|f_p^{(4)}(s)|\,ds=-f_p'''(c)=|f_p'''(c)|.
\]
Combining the explicit fourth-order term in \eqref{eq:EM} with its
remainder gives the factor $2/720=1/360$ in \eqref{eq:EM-bound}.
Substituting $h=1/t$ proves \eqref{eq:crossover}.

For the limit, apply the same formula two derivatives further. The
sixth-derivative remainder is $O(h^6)$ for fixed $c$, since the positive
inverse-power expansion gives integrability of every derivative from the
first onward. Thus the fourth-order coefficient is
$f_p'''(c)/720$ and the unrecorded terms are $O(h^6)$, proving
\eqref{eq:EM-limit}. Bounds on these derivatives are uniform on a compact
$c$-interval, which also proves the stated uniformity.
\end{proof}

\subsection{Reading the three regimes}
If $N/t\to\infty$ and $N\to\infty$, the elementary bounds already
imply
\begin{equation}\label{eq:large-N}
 S_{p,N}(t)\sim\frac{tN^{1-p}}{p-1}.
\end{equation}
Indeed, the ratio of the upper and lower bounds tends to one. When
$N/t\to c\in(0,\infty)$, the leading term is instead
$t^{2-p}\Phi_p(c)$. The exponent $p=2$ is critical: after factoring out
$x^2$ in the solution error \eqref{eq:error-exact}, the transition scale
has no extra power of $t$ at $p=2$, grows for $1<p<2$, and decays for
$p>2$.

At the critical exponent,
\begin{equation}\label{eq:Phi2}
 \Phi_2(c)=\int_0^{1/c}\frac{e^u-1}{u}\,du
          =\sum_{k=1}^{\infty}\frac{c^{-k}}{k\,k!}.
\end{equation}
For general $p>1$, the corresponding absolutely convergent series is
\begin{equation}\label{eq:Phi-series}
 \Phi_p(c)=\sum_{k=1}^{\infty}
   \frac{c^{-(p+k-2)}}{k!\,(p+k-2)}.
\end{equation}
These formulas also supply stable numerical evaluations of the transition
function.

The crossover is a joint truncation statement. It does not append a new
canonical term at exponent $2$ to the formal series, nor does it justify
exchanging an infinite Hahn sum with analytic evaluation. The dependence
on how $N$ grows with $t$ is exactly the information that a fixed finite
prefix leaves unresolved.

\subsection{Numerical illustration}
Table~\ref{tab:crossover} uses $p=2$, $c=1$, so $N=t$. Write
\[
 Q(t)=\Phi_2(1)-\frac{f_2(1)}{2t}-\frac{f'_2(1)}{12t^2}.
\]
The error column is computed from the convergent Hurwitz-zeta expansion
\begin{equation}\label{eq:hurwitz}
 S_{p,N}(t)=\sum_{k=1}^{\infty}
                  \frac{t^k}{k!}\zeta(p+k-1,N+1),
\end{equation}
which follows by absolute summation exactly as in
\eqref{eq:zeta-expansion}. The bound is a theorem, while the displayed
decimals are high-precision floating-point illustrations, not interval
certificates.

\begin{table}[htbp]
\centering
\small
\begin{tabular}{rrrr}
\toprule
$t=N$ & $S_{2,N}(t)$ & $S_{2,N}(t)-Q(t)$ & Proved error bound\\
\midrule
20 & 1.275868641749 & $-7.48089\times10^{-7}$ & $1.50037\times10^{-6}$\\
40 & 1.296654652769 & $-4.68538\times10^{-8}$ & $9.37734\times10^{-8}$\\
80 & 1.307220654853 & $-2.92990\times10^{-9}$ & $5.86084\times10^{-9}$\\
160 & 1.312546962496 & $-1.83143\times10^{-10}$ & $3.66302\times10^{-10}$\\
\bottomrule
\end{tabular}
\caption{The critical accumulation crossover. Decimal computations use
85-digit working precision; the bound is \eqref{eq:EM-bound}.}
\label{tab:crossover}
\end{table}

\section{Computation, verification, and a route to Lean}
\label{sec:verification}

\subsection{Finite exact normalizer}
The supplied \path{verification/verify.py} represents a matrix series as a
dictionary from rational exponents to exact SymPy matrices. For a finite
positive rational support and a rational cutoff, it enumerates the generated
semigroup below the cutoff, processes exponents in increasing order, and
uses \eqref{eq:T} and \eqref{eq:recurrence}. After constructing $H$ and
$R$, it checks the complete coefficient residual of \eqref{eq:gauge},
including all matrix products, at every enumerated exponent. It also
checks the normalization masks and nilpotence of $B$.

The code accepts commuting nilpotent $N$ rather than assuming a semisimple
constant matrix. This matters: the finite geometric inverse and its
noncommutative products are among the places most susceptible to a sign
or projection error. Exact checks include the symbolic polynomial
$e+c(a-d)$, sharp chains of ranks $1$ through $7$, and examples with a
nontrivial repeated-eigenvalue Jordan block.

The recorded run, with seed $20260929$, has the following scope:
\begin{center}
\begin{tabular}{lr}
\toprule
Verification item & Count\\
\midrule
Exact finite gauge systems & 33\\
Exact coefficient checks, including scalar counterexample checks & 292\\
Crossover parameter cases & 36\\
\bottomrule
\end{tabular}
\end{center}
The numerical cases use $p\in\{1.5,2,2.5\}$,
$c\in\{0.5,1,2\}$, and $t\in\{20,40,80,160\}$. They compare
\eqref{eq:hurwitz} with both \eqref{eq:tail-bounds} and
\eqref{eq:EM-bound}, and record convergence toward
\eqref{eq:EM-limit}. The environment was Python~3.13.5,
SymPy~1.14.0, and mpmath~1.3.0.

These are reproducible finite consistency checks, not substitutes for the
proofs. They do not establish the infinite-support recursion, the universal
quantifiers in Theorem~\ref{thm:universal}, or arbitrary-support
well-ordering. No Lean code in this package has been claimed to prove those
statements.

\subsection{A formalization decomposition}
A direct Lean development can be organized around four interfaces, each
with a narrow mathematical contract.

\paragraph{Support arithmetic.}
Reuse the repository's two-factor additive Neumann bridge to prove the
positive-semigroup finite-word lemma. The additional content is the real
lower bound $\delta$ and the conversion of a fixed exponent bound into a
natural-number word-length bound. Do not assume that bounded supports are
finite.

\paragraph{Finite-dimensional homological algebra.}
Formalize the spectral-difference masks, their commutation with $\ad N$,
the identity $(\ad N)^{2m-1}=0$, and the finite inverse
\eqref{eq:T}. This component is independent of Hahn series and can be
checked over a characteristic-zero field with suitable eigenvalue data.

\paragraph{Strong Hahn derivation and recursive gauge.}
Lift those operators coefficientwise and prove the valuation improvement
\eqref{eq:valuation-contraction}. The derivative's product rule, the
normal-form equation, and constants in $\K[L]$ should be explicit
interfaces, not informal consequences of a field name. The repository's
block-antiderivative API is related algebra, but it is not already a
formalization of this matrix construction.

\paragraph{Analytic realization.}
Build the weighted $\ell^1$ convolution algebra, prove the operator norm
bounds and contraction estimates, and only then supply a branchwise
analytic evaluation map. The universal converse can be formalized
separately with the rank-one family. This prevents analytic summation from
being silently built into the formal Hahn datatype.

The suggested theorem targets are therefore finite-word support,
normalized gauge existence and uniqueness, nilpotent log filtration,
finite ancestry, absolute-convergence preservation, and its sharp failure
criterion. None is asserted to be present in the inspected repository
merely because related names already exist there.

\section{Further research questions and concrete next targets}
\label{sec:future}

\subsection{Universal convergence for a fixed smaller input support}
Theorem~\ref{thm:universal} quantifies over all inputs supported in $\G$.
Fix instead a generating set $E\subsetneq\G$ and allow coefficients only
on $E$. Is the absence of a left-accumulating spectral difference in $\G$
still necessary? A bad exponent might be reachable only through a long
product whose coefficients or matrix directions are constrained. A useful
first target is a triangular three-dimensional system with two permitted
matrix units and an accumulating two-letter sumset. This asks whether an
additive obstruction must have a realizable noncommutative path.

\subsection{Weighted convergence in the zero-gap case}
For the square-zero family, convergence is exactly a weighted summability
condition after division by $\gamma-\rho$. For general matrices, products
of several small divisors appear. Construct a weight on exponent paths that
is both sufficient for a Banach fixed-point argument and close to necessary.
One candidate is the least submultiplicative majorant of the homological
inverse norms. The main issue is whether cancellations permit substantially
larger classes than any absolute path majorant can capture.
% ed. (2026-09-29, batch 52): the two questions above are answered, for a restricted class, by a later package; recorded on filing.
\begin{ednote}
The two preceding questions (a fixed smaller input support; weighted
convergence in the zero-gap case) are answered, for finite acyclic patterns
with a diagonal real constant term, a well-ordered support on each edge and
independently variable edge coefficients, by the later package
\path{Analysis/Transseries/docs/series-and-transseries/Path_Sensitive_Small_Divisors_Hahn_Fuchsian/path_sensitive_hahn.tex}
(batch~52; not reviewed). Universal convergence holds exactly when no
endpoint spectral difference of a directed path is a left accumulation
point of that path's own exponent sumset (its \texttt{thm:gap}), so an
accumulating difference in \(\G\) that no path realizes is harmless; for
independent weighted inputs it holds exactly when every path-kernel norm
is finite, and that norm is the exact multilinear operator norm of the
path contribution (its \texttt{thm:weighted}). Correlated inputs can do
better: its example in Section~\texttt{sec:cancellation} cancels two
accumulating paths. Its rank-\(r\) resonant chain converges exactly when
\(\sum n^{r-1}\abs{c_n}<\infty\) (its \texttt{thm:chain}), which at
\(r=2\) is the threshold \(p>2\) of Proposition~\ref{prop:threshold}. For
the complete triangular pattern with the whole semigroup on every edge, its
\texttt{thm:gap} reduces to the triangular case of
Theorem~\ref{thm:universal} with a diagonal constant term, consistently;
Theorem~\ref{thm:universal} also allows nontriangular perturbations and a
nilpotent part, which that package excludes.
Nonlinear counterparts of Theorem~\ref{thm:universal} are proved in
\path{Nonlinear_Hahn_Dulac_Finite_Resonance_Control/} (its
\texttt{thm:universal}) and
\path{Nonlinear_Hahn_Fuchsian_Algebraic_Convergence_Loci/} (its
\texttt{thm:gap}), in the same directory.
\end{ednote}

\subsection{Renormalization compatible with multiplication and gauge change}
The difference sum in \eqref{eq:Psi} gives one natural analytic realization
of a divergent normalizer, selected by a boundary value at $x=1$. Can this
be extended to a summation operation on an algebra of accumulating-support
series that commutes with $D$ and respects products? Different normalizing
conditions add resonant homogeneous terms. A precise first task is to
classify all such realizations for upper-triangular rank-three systems and
identify the additional compatibility identities caused by nonzero
products of matrix units.

\subsection{Uniform spectral parameters and moving accumulation resonances}
Let the positive eigenvalue difference vary through an accumulation point
of $\G$. The constant $C$ in \eqref{eq:C} blows up, even though individual
solutions may continue after renormalization. Determine a uniform scaling
limit with $\rho-\gamma_n$ and the spectral displacement competing. The
model $\gamma_n=2-1/n$ suggests transition parameters involving both
$N/\log(1/x)$ and the displacement times $N$. A uniform theorem should
separate exact resonance, near-resonance, and the accumulation boundary.

\subsection{Iterated-logarithmic coefficient fields}
Replace the constant matrices $A_\alpha$ by controlled Hahn series in
$\log(1/x)$ and further iterated logarithms. Which parts of the finite
matrix certificate survive, and when does one additional logarithmic level
suffice? This connects directly with the earlier residue-depth results in
ProveIt. A first tractable case keeps the outer exponent gap positive and
allows one rational function of $\log(1/x)$ in each coefficient. The
scalar block inverse then interacts with the matrix spectral inverse.

\subsection{Higher-rank exponent groups and the loss of finite word depth}
In a non-Archimedean exponent group, a positive generator can be smaller
than every positive real multiple of another scale. The bound
$k\delta\leq r_*$ need not bound $k$. Find an intrinsic replacement for
the finite-ancestry hypothesis, perhaps a no-infinite-resonant-path
condition. The intended conclusion is not necessarily a finite ordinary
matrix but a controlled transfinite filtration. A counterexample to a naive
finite-depth extension would be as informative as a sufficient theorem.

\subsection{Effective critical ancestry and exact real presentations}
For finite rational generators, the included algorithm is elementary.
For finite algebraic generators, exact comparison and equality are still
available but implementation costs differ. For infinite computable supports,
identify representation hypotheses under which $E_{\mathrm{crit}}$ can
actually be extracted. A useful deliverable would be a theorem separating
finite dependence, semidecidable resonance detection, and decidable complete
certificates, rather than conflating these three levels.

\subsection{Geometry of the logarithmic strata}
The loci $B^{d+1}=0$ are algebraic in critical coefficient data. Determine
their irreducible components and generic codimensions for a prescribed
spectral graph. Even a three- or four-level chain can have cancellation
between several paths. The polynomial $e+c(a-d)$ is the first nontrivial
example. Higher-rank cases could give a useful classification of when
logarithms disappear under parameter tuning and when a small perturbation
necessarily restores them.

\subsection{Irregular systems and exponential transseries sectors}
For an irregular system, exponentially small or large normal-form sectors
introduce a second source of divergence, distinct from accumulation at a
finite spectral difference. Develop a model in which one irregular
exponential sector and one accumulating real-exponent support coexist.
The immediate goal is a theorem separating a Borel-summation obstruction
from the absolute small-divisor obstruction proved here. Existing
resurgence terminology alone does not identify which mechanism is active.

\subsection{Formal verification of the sharp equivalence}
The most valuable initial proof-assistant target is not the numerical
crossover but Theorem~\ref{thm:universal}: its quantifiers expose the
formal/analytic distinction particularly clearly. A staged project can
first verify the square-zero counterexample and the finite-dimensional
inverse, then the $\ell^1$ contraction, and finally the positive-Hahn
support layer. The finite-ancestry certificate would make a natural exact
computational interface on top of that verified core.

\section{Conclusion}
\label{sec:conclusion}

The formal normal form and the analytic convergence problem have different
answers. Positive well-ordered supports suffice for a unique normalized
Hahn gauge. Exact resonances reduce to a finite nilpotent matrix, and even
an infinitely populated low-exponent region contributes only finitely many
ancestors to its logarithmic obstruction. But absolute convergence depends
on an additional relation between the exponent semigroup and the spectrum.
The required relation is exactly a uniform gap from the nonresonant
spectral differences.

The explicit family with exponents $2-1/n$ demonstrates all the distinctions
at once. Its coefficients converge for $p>1$, its formal normalizer
converges only for $p>2$, and its renormalized analytic solution realizes
every finite asymptotic prefix throughout the whole range $p>1$. The
transition $N\asymp\log(1/x)$ quantifies the information lost when an
accumulating formal expansion is treated as an ordinary convergent sum.

These results offer a specific extension of the repository's support,
resonance, and flatness calculus. They also delimit that extension: the
proofs concern the stated real-exponent Fuchsian class, not general
transseries equations; the certificate is mathematically finite but not
automatically computable from arbitrary infinite presentations; and
historical priority remains a separate question from the correctness of
the supplied proofs.

\appendix
\section{Proof and provenance audit}
\label{sec:audit}

\subsection{Assumptions that must not be dropped}
The formal construction uses positive real input exponents with a common
positive least element, well ordering, finite matrix dimension, and a
real-spectrum constant matrix. The log-degree classification fixes the
ambient algebra $\K[L]$. The analytic criterion uses absolute
$\ell^1$ convergence in a fixed Jordan basis and quantifies over the whole
positive semigroup. The uniform crossover holds for $p>1$ and for $N\geq1$;
its compact-parameter assertion keeps $N/t$ away from zero and infinity.

Each restriction has a role. The positive lower bound controls the number
of input factors. Well ordering controls finite coefficient convolution.
Real eigenvalues allow real monomials $x^{\lambda_i}$ in the chosen Hahn
field. Absolute norms justify products and analytic evaluation. The
whole-semigroup quantifier is needed for the converse construction. The
fixed logarithmic algebra makes minimal log degree a well-posed invariant.

\subsection{Dependency map}
Lemma~\ref{lem:words} and Lemma~\ref{lem:inverse} imply the formal
normalization. The normal form, the constant-field lemma, and the
nilpotence of $B$ imply the log filtration. Finite-word dependence gives
the ancestry certificate. The analytic forward direction adds a uniform
operator bound and Banach contraction; the converse uses only the spectral
gap characterization and a rank-one eigenmatrix. The model and crossover
are proved independently by scalar summation and Euler--Maclaurin, so they
do not depend on unproved analytic realization of arbitrary Hahn series.

\subsection{Repository inspection boundary}
The snapshot was pinned to \snap. The group index, canonical-package
status documentation, the two-factor Hahn-support Lean module, the
block-antiderivative module, and the earlier residue-obstruction package
README were inspected through the connected GitHub tools. The canonical
large TeX volume exceeded the connector's content limit; no claim of an
exhaustive line-by-line audit of that volume or the whole repository is made.
The supplied provenance JSON records file paths, blob identifiers, and the
scope of inspected ranges. No repository file was modified.

\subsection{Novelty and verification boundary}
The normal-form architecture specializes to classical Fuchsian theory, and
the finite-obstruction philosophy has an explicit repository antecedent.
The spectral accumulation equivalence, its sharp analytic examples, the
finite-ancestry formulation, and the uniform crossover are the proposed
research additions developed here. The literature check was targeted, not
exhaustive; no independent peer review has occurred. The exact scripts
check finite algebraic instances, while the numerical results illustrate
bounds proved in the text. Neither is represented as a machine proof of
the general theorems.

\begin{thebibliography}{99}
\bibitem{repo-index}
Vladimir Reshetnikov, \emph{ProveIt: Series and transseries}, repository
group index, snapshot \snap, inspected 29 September 2026.
\url{https://github.com/VladimirReshetnikov/ProveIt/blob/04473354a0f3edff2d6c365caf26170ca6b88735/Analysis/Transseries/docs/series-and-transseries/README.md}.

\bibitem{repo-wellbased}
Vladimir Reshetnikov, ProveIt contributors,
\emph{TransseriesWellBased.lean}, same repository snapshot.
\url{https://github.com/VladimirReshetnikov/ProveIt/blob/04473354a0f3edff2d6c365caf26170ca6b88735/Analysis/Transseries/Lean/Transseries/TransseriesWellBased.lean}.
The additive two-factor interfaces used for the crosswalk are
\texttt{Fabius.neumann\_add\_isPWO} and
\texttt{Fabius.neumann\_finite\_decompositions}.

\bibitem{repo-block}
Vladimir Reshetnikov, ProveIt contributors,
\emph{TransseriesBlockAntiderivative.lean}, same repository snapshot.
\url{https://github.com/VladimirReshetnikov/ProveIt/blob/04473354a0f3edff2d6c365caf26170ca6b88735/Analysis/Transseries/Lean/Transseries/TransseriesBlockAntiderivative.lean}.

\bibitem{repo-residue}
\emph{Finite Residue Obstructions and Logarithmic-Depth Promotion in Hahn
Transseries}, ProveIt research-package README, dated 29 September 2026,
same repository snapshot.
\url{https://github.com/VladimirReshetnikov/ProveIt/blob/04473354a0f3edff2d6c365caf26170ca6b88735/Analysis/Transseries/docs/series-and-transseries/Residue_Obstructions_Logarithmic_Depth_Promotion/README.md}.
The comparison in this article is with the results and scope stated in
that package's README.

\bibitem{tanny-yakovenko}
Shira Tanny and Sergei Yakovenko,
\emph{On local Weyl equivalence of higher order Fucshian equations},
arXiv:1412.7830 (2014).
\url{https://arxiv.org/abs/1412.7830}.
The spelling of ``Fucshian'' follows the arXiv title; p.~6 of the PDF
records the classical resonant polynomial normal form.

\bibitem{vdh}
Joris van der Hoeven,
\emph{Operators on generalized power series},
Illinois Journal of Mathematics \textbf{45} (2001), no.~4,
DOI: 10.1215/ijm/1258138061.
\url{https://www.texmacs.org/joris/noeth/noeth.html}.

\bibitem{joswig-smith}
Michael Joswig and Ben Smith,
\emph{Convergent Hahn Series and Tropical Geometry of Higher Rank},
arXiv:1809.01457.
\url{https://arxiv.org/abs/1809.01457}.
Cited for the broader study of convergent Hahn series, not as a source for
the differential-equation criterion proved here.
\end{thebibliography}
\end{document}
